% Image credits for the photographs and AI illustrations of Book 2
% (University Chemistry, Year 1). Machine-readable license records live in
% images/book2/CREDITS.md; keep the two files in sync. The Book 2 writer
% owns this file: add one \item per photograph (in an itemize, after the
% first paragraph) as photographs land.
% Arabic edition: personal names are transliterated as in the captions, with
% the original spelling in parentheses (the legally required credit); user
% names, institutions, works and licences stay verbatim (as in the Arabic
% Book 1 credits page, image-credits.ar.tex).
\chapter*{حقوق الصور}

رسوم هذا الكتاب أصلية. أما الصور الفوتوغرافية، حيث استُعملت، فمستعملة
بمقتضى رخصها الحرة، مع الشكر لأصحابها؛ وروابط المصادر الكاملة وروابط
الرخص لكل صورة مدرجة في
\texttt{images/book2/CREDITS.md} في مستودع الكتاب.
والرسوم التحذيرية للأخطار هي رسوم النظام المنسَّق عالميًا لتصنيف
المواد الكيميائية ووسمها، كما رُسمت لحزمة
\texttt{ghsystem} في \TeX\ Live.

\bigskip
وإلى جانب الصور الفوتوغرافية والرسوم الأصلية، بعض الأشكال
رسوم مولَّدة بالذكاء الاصطناعي، أُنتجت بأداة توليد الصور من OpenAI انطلاقًا من
أوامر كتبها محررو الكتاب، ورُوجعت من حيث الدقة الكيميائية قبل
إدراجها. والأمر الكامل لكل صورة مولَّدة
مدوَّن في \texttt{images/book2/ai/PROMPTS.md} في المستودع.

\bigskip
\noindent الصور الفوتوغرافية، حسب الفصل:
\begin{itemize}\small
  \item الفصل~1: نيلز بور (Niels Bohr)، نحو 1922، AB Lagrelius \& Westphal، ملك عام
        (Wikimedia Commons).
  \item الفصل~2: لينوس بولينغ (Linus Pauling)، 1962، Nobel Foundation، ملك عام
        (Wikimedia Commons).
  \item الفصل~5: عيّنة نحاس طبيعي، التقطها Flipin، OG، CC0 (Wikimedia
        Commons).
  \item الفصل~6: بلورات الهاليت، والفلوريت وماس ثماني الأوجه، التقطها
        جيمس سانت جون (James St. John)، CC BY 2.0؛ والغرافيت، التقطها زبينيك بوريفال
        (Zbynek Burival)، CC BY-SA 4.0؛ وبلورات الثلج التقطها ويلسون أ. بنتلي
        (Wilson A. Bentley) سنة 1902، ملك عام (كلها من Wikimedia Commons).
  \item الفصل~8: سفانتي أرهينيوس (Svante Arrhenius)، 1909، Meisenbach Riffarth \& Co.، ملك
        عام (Wikimedia Commons).
  \item الفصل~11: رواسب كلوريد الفضة وبروميدها ويوديدها، التقطها
        Rrausch1974، CC BY-SA 3.0 (Wikimedia Commons).
  \item الفصل~12: معلَّق هيدروكسيد النحاس(II)، التقطه Alvy16، CC BY 4.0؛
        ومحلول أمين النحاس، التقطه Chemicalinterest، ملك عام (كلاهما من
        Wikimedia Commons).
  \item الفصل~13: فالتر نرنست (Walther Nernst)، 1889، Smithsonian Libraries، ملك عام
        (Wikimedia Commons).
  \item الفصل~16: ياكوبوس هنريكوس فانت هوف (Jacobus Henricus van 't Hoff)، صورة نُشرت سنة 1911،
        ملك عام (Wikimedia Commons).
  \item الفصل~21: فيكتور غرينيار (Victor Grignard)، ملك عام (Wikimedia Commons).
  \item الفصل~25: فلاديمير ماركوفنيكوف (Vladimir Markovnikov)، صورة من نعي
        نُشر سنة 1905، ملك عام (Wikimedia Commons).
  \item الفصل~26: أحواض الليثيوم في سبخة أتاكاما، NASA Earth
        Observatory، ملك عام؛ والصوديوم في الزيت المعدني، التقطها جاستن أورجيتيس (Justin Urgitis)،
        CC BY-SA 2.5؛ وألوان اللهب، التقطها Hegelrast، CC BY-SA 4.0 (كلها من
        Wikimedia Commons).
  \item الفصل~27: بوكسيت حمّصي، التقطه جيمس سانت جون (James St. John)، CC BY 2.0؛ والبوراكس، التقطه
        آرام دوليان (Aram Dulyan)، ملك عام؛ والكوارتز، التقطته ستيفاني كليفورد (Stephanie Clifford)، CC BY 2.0؛
        والفوسفور الأبيض، Hi-Res Images of Chemical Elements، CC BY 3.0؛
        وفريتس هابر (Fritz Haber)، The Nobel Foundation، ملك عام (كلها من Wikimedia
        Commons).
  \item الفصل~28: استخراج الكبريت في فوهة كاوا إيجين، التقطها Sémhur،
        CC BY-SA 4.0؛ والكلور في أمبولة، التقطها ف.~أولن (W.~Oelen)، CC BY-SA 3.0؛
        والبروم في أمبولة، التقطها Jurii، CC BY 3.0؛ وبلورات اليود، التقطها
        Dnn87، CC BY 3.0 (كلها من Wikimedia Commons).
  \item الفصل~29: منصة ساخنة (منصة كوفلر)، التقطها HBR، CC BY-SA 3.0؛ ومقياس انكسار
        آبي، التقطه Foreade، CC BY-SA 4.0 (كلاهما من Wikimedia Commons).
\end{itemize}
\clearpage
